\section{Related Work}
\subsection{Automatic Lyric Transcription (ALT)}
Automatic Lyric Transcription aims to decode text from polyphonic music recordings or isolated vocal tracks \cite{mesaros2010automatic,gupta2022deep}. Mesaros and Virtanen \cite{mesaros2010automatic} identified three fundamental acoustic hurdles distinguishing singing from speech: prolonged vowel durations, pitch vibrato modulations, and formant shifts induced by expressive vocal production. Recent studies have investigated self-supervised speech representations such as wav2vec 2.0 \cite{baevski2020wav2vec} and HuBERT \cite{hsu2021hubert}. Ou et al. \cite{ou2022transfer} demonstrated that fine-tuning wav2vec 2.0 yields promising results on English lyrics but struggles with sustained notes in the absence of explicit temporal guidance. 

Large-scale generative speech foundation models, notably OpenAI Whisper \cite{radford2023robust}, exhibit remarkable robustness on spoken audio but frequently break down on musical singing. Zhuo et al. \cite{zhuo2023lyricwhiz} proposed LyricWhiz, an unconstrained zero-shot framework combining Whisper transcription with large language model post-processing. Stoller et al. \cite{stoller2024jamalt} introduced the Jam-ALT benchmark, revealing that conventional Word Error Rate (WER) fails to account for human reading perception during musical alignment. In classical Western repertoire, Berendes et al. \cite{berendes2024lyrics} analyzed Whisper on Schubert's \textit{Winterreise}, demonstrating that model degradation stems primarily from expressive vocalization dynamics rather than backing instrumental accompaniment. To separate vocal stems from polyphonic mixtures, state-of-the-art hybrid separation architectures such as Demucs \cite{defossez2021hybrid} and HTDemucs \cite{rouard2023hybrid} provide essential acoustic isolation.

\subsection{Tonal Languages and Tone-Melody Mismatch}
In tonal phonology, pitch trajectory is lexically distinctive. For Vietnamese, each syllable carries one of six canonical tones that dictate meaning \cite{pham2003vietnamese,nguyen1997vietnamese,brunelle2009tone}. In music composition, Kirby and Ladd \cite{kirby2016tone} investigated tone-melody correspondence in Vietnamese popular song, proving that while songwriters generally align musical pitch contours with lexical tones, musical melodic constraints frequently force tone flattening or inverted directionality. Wong and Diehl \cite{wong2002how} observed that the auditory processing of lexical tone and musical pitch interacts closely in native tonal speakers.

Recent developments in Vietnamese speech technology have produced large-scale corpora, including VIVOS \cite{luong2016vivos}, Mozilla Common Voice \cite{ardila2020common}, VietSuperSpeech \cite{thanh2026vietsuperspeech}, and VietPrism \cite{hoang2026vietprism}. Furthermore, dedicated singing corpora such as VietSing \cite{vu2024vietsing} have emerged. However, systematic empirical investigations addressing the intersection of tonal phonology, cross-attention mechanics, and parameter-efficient model adaptation in Vietnamese singing ASR remain distinctly scarce.
